% =============================================================================
% Conclusion — conservative, evidence-aligned framing
% Drop-in include: % =============================================================================
% Conclusion — conservative, evidence-aligned framing
% Drop-in include: % =============================================================================
% Conclusion — conservative, evidence-aligned framing
% Drop-in include: % =============================================================================
% Conclusion — conservative, evidence-aligned framing
% Drop-in include: \input{sections/conclusion}
% Target: <= 0.75 page
% =============================================================================
\section{Conclusion}
\label{sec:conclusion}

\ours{} provides a shared empirical substrate for studying RL post-training
across algorithms, frameworks, and model families, but the evidence it offers
is narrower than a broad ``five laws'' framing. The current release most
strongly supports three conservative claims. First, ZVF/GU are useful
descriptive diagnostics of when binary-reward GRPO stops producing informative
within-group signal; they should not yet be read as standalone causal or
incrementally predictive statistics. Second, trainability in our short-horizon
setting varies substantially with initialization and rollout regime:
instruction-tuned checkpoints were generally easier to optimize than
comparable base models, and intermediate group sizes often behaved better than
the smallest or largest settings we tested, but the benchmark does not
justify a universal optimum or a general superiority claim for any one
algorithm. Third, PPO-vs.-GRPO rankings and frontier-run stability are
heterogeneous across model families and stacks, so our results argue against
one-size-fits-all recommendations rather than for a new one.

The most important negative result is also the most useful one: on held-out
GSM8K, the mean Qwen3-8B GRPO gain over the base model is small and not
statistically significant under our current evaluation. Tool-use and code
results are even less conclusive because they rely on sparse or custom reward
protocols. That boundary on what we can claim is not a weakness to hide; it is
the point of releasing the benchmark.

The immediate next steps are experimental rather than rhetorical: multivariate
tests of ZVF against reward mean, entropy, and divergence proxies; token-budget
matched multi-seed PPO/GRPO comparisons in a single open stack; broader
held-out evaluation without checkpoint cherry-picking; and non-math tasks with
process rewards or richer execution-based metrics. The released traces,
checkpoints, and scripts are intended to make that stricter follow-up easy.

% Target: <= 0.75 page
% =============================================================================
\section{Conclusion}
\label{sec:conclusion}

\ours{} provides a shared empirical substrate for studying RL post-training
across algorithms, frameworks, and model families, but the evidence it offers
is narrower than a broad ``five laws'' framing. The current release most
strongly supports three conservative claims. First, ZVF/GU are useful
descriptive diagnostics of when binary-reward GRPO stops producing informative
within-group signal; they should not yet be read as standalone causal or
incrementally predictive statistics. Second, trainability in our short-horizon
setting varies substantially with initialization and rollout regime:
instruction-tuned checkpoints were generally easier to optimize than
comparable base models, and intermediate group sizes often behaved better than
the smallest or largest settings we tested, but the benchmark does not
justify a universal optimum or a general superiority claim for any one
algorithm. Third, PPO-vs.-GRPO rankings and frontier-run stability are
heterogeneous across model families and stacks, so our results argue against
one-size-fits-all recommendations rather than for a new one.

The most important negative result is also the most useful one: on held-out
GSM8K, the mean Qwen3-8B GRPO gain over the base model is small and not
statistically significant under our current evaluation. Tool-use and code
results are even less conclusive because they rely on sparse or custom reward
protocols. That boundary on what we can claim is not a weakness to hide; it is
the point of releasing the benchmark.

The immediate next steps are experimental rather than rhetorical: multivariate
tests of ZVF against reward mean, entropy, and divergence proxies; token-budget
matched multi-seed PPO/GRPO comparisons in a single open stack; broader
held-out evaluation without checkpoint cherry-picking; and non-math tasks with
process rewards or richer execution-based metrics. The released traces,
checkpoints, and scripts are intended to make that stricter follow-up easy.

% Target: <= 0.75 page
% =============================================================================
\section{Conclusion}
\label{sec:conclusion}

\ours{} provides a shared empirical substrate for studying RL post-training
across algorithms, frameworks, and model families, but the evidence it offers
is narrower than a broad ``five laws'' framing. The current release most
strongly supports three conservative claims. First, ZVF/GU are useful
descriptive diagnostics of when binary-reward GRPO stops producing informative
within-group signal; they should not yet be read as standalone causal or
incrementally predictive statistics. Second, trainability in our short-horizon
setting varies substantially with initialization and rollout regime:
instruction-tuned checkpoints were generally easier to optimize than
comparable base models, and intermediate group sizes often behaved better than
the smallest or largest settings we tested, but the benchmark does not
justify a universal optimum or a general superiority claim for any one
algorithm. Third, PPO-vs.-GRPO rankings and frontier-run stability are
heterogeneous across model families and stacks, so our results argue against
one-size-fits-all recommendations rather than for a new one.

The most important negative result is also the most useful one: on held-out
GSM8K, the mean Qwen3-8B GRPO gain over the base model is small and not
statistically significant under our current evaluation. Tool-use and code
results are even less conclusive because they rely on sparse or custom reward
protocols. That boundary on what we can claim is not a weakness to hide; it is
the point of releasing the benchmark.

The immediate next steps are experimental rather than rhetorical: multivariate
tests of ZVF against reward mean, entropy, and divergence proxies; token-budget
matched multi-seed PPO/GRPO comparisons in a single open stack; broader
held-out evaluation without checkpoint cherry-picking; and non-math tasks with
process rewards or richer execution-based metrics. The released traces,
checkpoints, and scripts are intended to make that stricter follow-up easy.

% Target: <= 0.75 page
% =============================================================================
\section{Conclusion}
\label{sec:conclusion}

\ours{} provides a shared empirical substrate for studying RL post-training
across algorithms, frameworks, and model families, but the evidence it offers
is narrower than a broad ``five laws'' framing. The current release most
strongly supports three conservative claims. First, ZVF/GU are useful
descriptive diagnostics of when binary-reward GRPO stops producing informative
within-group signal; they should not yet be read as standalone causal or
incrementally predictive statistics. Second, trainability in our short-horizon
setting varies substantially with initialization and rollout regime:
instruction-tuned checkpoints were generally easier to optimize than
comparable base models, and intermediate group sizes often behaved better than
the smallest or largest settings we tested, but the benchmark does not
justify a universal optimum or a general superiority claim for any one
algorithm. Third, PPO-vs.-GRPO rankings and frontier-run stability are
heterogeneous across model families and stacks, so our results argue against
one-size-fits-all recommendations rather than for a new one.

The most important negative result is also the most useful one: on held-out
GSM8K, the mean Qwen3-8B GRPO gain over the base model is small and not
statistically significant under our current evaluation. Tool-use and code
results are even less conclusive because they rely on sparse or custom reward
protocols. That boundary on what we can claim is not a weakness to hide; it is
the point of releasing the benchmark.

The immediate next steps are experimental rather than rhetorical: multivariate
tests of ZVF against reward mean, entropy, and divergence proxies; token-budget
matched multi-seed PPO/GRPO comparisons in a single open stack; broader
held-out evaluation without checkpoint cherry-picking; and non-math tasks with
process rewards or richer execution-based metrics. The released traces,
checkpoints, and scripts are intended to make that stricter follow-up easy.
